\DocumentMetadata{
  pdfversion=2.0,
  pdfstandard=ua-2,
  lang=en-US,
  tagging=on,
  tagging-setup={math/setup=mathml-SE,tabsorder=structure}
}
\documentclass[11pt,letterpaper]{article}
\usepackage[margin=0.68in,includefoot]{geometry}
\usepackage{newcomputermodern}
\usepackage{amsmath,amscd,mathtools}
\usepackage{tikz,adjustbox}
\providecommand{\square}{\mdlgwhtsquare}
\usepackage[hidelinks]{hyperref}
\hypersetup{
  pdftitle={Algebra First-Year Exam, August 2017, Part 2},
  pdfauthor={University of Florida Department of Mathematics},
  pdfsubject={Graduate mathematics examination},
  pdfdisplaydoctitle=true,
  bookmarksopen=true
}
\setcounter{secnumdepth}{0}
\setlength{\parindent}{0pt}
\setlength{\parskip}{0.28em}
\setlength{\emergencystretch}{3em}
\setlength{\leftmargini}{2.15em}
\setlength{\leftmarginii}{2.65em}
\setlength{\leftmarginiii}{3em}
\renewcommand{\labelenumi}{\textbf{\theenumi.}}
\pagestyle{empty}
\raggedbottom
\allowdisplaybreaks
\newenvironment{examproblems}[1]{%
  \begin{enumerate}%
  \setcounter{enumi}{#1}%
  \setlength{\topsep}{0.35em}%
  \setlength{\partopsep}{0pt}%
  \setlength{\itemsep}{0.48em}%
  \setlength{\parsep}{0.18em}%
}{\end{enumerate}}
\newenvironment{examsubparts}{%
  \begin{enumerate}%
  \setlength{\topsep}{0.18em}%
  \setlength{\partopsep}{0pt}%
  \setlength{\itemsep}{0.16em}%
  \setlength{\parsep}{0.1em}%
}{\end{enumerate}}
\newenvironment{examinstructions}{%
  \par\begingroup\itshape
}{%
  \par\endgroup\vspace{0.18em}
}
\makeatletter
\renewcommand\section{\@startsection{section}{1}{0pt}%
  {-1.25ex plus -.25ex minus -.1ex}%
  {0.55ex plus .1ex}%
  {\normalfont\large\bfseries}}
\renewcommand{\@maketitle}{%
  \newpage\null\vskip -1em%
  {\footnotesize\bfseries
    \href{https://www.ufl.edu/}{University of Florida}, \href{https://math.ufl.edu/}{Department of Mathematics}\par}%
  \vskip -0.5em%
  \tagstructbegin{tag=Title}%
    {\LARGE\bfseries\href{https://gma.math.ufl.edu/exams/algebra-first-year/}{Algebra First-Year Exam}, August 2017, Part 2\par}%
  \tagstructend%
  \vskip 0.25em%
  \tagmcbegin{artifact}%
    \hrule height 1pt%
  \tagmcend%
  \vskip 1em%
}
\makeatother
\title{Algebra First-Year Exam, August 2017, Part 2}
\author{}
\date{}
\begin{document}
\maketitle
\thispagestyle{empty}
\begin{examinstructions}
Answer four problems. Write your answers clearly in complete English sentences. You may quote results (within reason) as long as you state them clearly.
\end{examinstructions}
\begin{examproblems}{0}
\item
Show that the ring \(\mathbb{Z}[\sqrt{-5}]\) is not a principal ideal domain, and give an example of a non-principal ideal in this ring.
\item
Let~\(R\) be a commutative ring with identity. Show that an element \(a \in R\) is contained in every maximal ideal if and only if \(1+ab\) is a unit in~\(R\) for all \(b \in R\).
\item
Show that if~\(p\) is an odd prime then the polynomial \(\frac{(x+2)^p-2^p}{x}\) is irreducible in \(\mathbb{Q}[X]\).
\item
Let~\(K\) be a field. Show that any~\(K\)-vector space~\(V\) has a basis (do not assume that~\(V\) has finite dimension).
\item
Suppose~\(R\) is a principal ideal domain and~\(M\) is a finitely generated~\(R\)-module. For any prime \(p \in R\), let \(M_p\) be the set of \(m \in M\) such that \(p^n m=0\) for some~\(n\).
\begin{examsubparts}
\item[\textnormal{(i)}]
Show that \(M_p\) is a submodule of~\(M\).
\item[\textnormal{(ii)}]
Show that if \(rM=0\) for some nonzero \(r \in R\), then~\(M\) is isomorphic to the direct sum of the \(M_p\) for all primes~\(p\) dividing~\(r\).
\end{examsubparts}
\item
Find representatives of all conjugacy classes of elements of order~\(5\) in \(GL_2(\mathbb{F}_{19})\). Hint: 
\[
x^5-1=(x-1)(x^2-4x+1)(x^2+5x+1)
\]
 in \(\mathbb{F}_{19}[x]\), and the quadratic factors are irreducible.
\end{examproblems}
\end{document}
